\documentclass[11pt]{ctexart}
\usepackage[a4paper,margin=1in]{geometry}
\usepackage{longtable,booktabs}
\title{Geant4-Agent 回归测试报告（中文）}
\author{自动化评测流程}
\date{2026-03-07\_232752}
\begin{document}
\maketitle
\section{统一 Pass/Fail 总表}
\begin{longtable}{p{0.06\linewidth}p{0.08\linewidth}p{0.07\linewidth}p{0.08\linewidth}p{0.07\linewidth}p{0.07\linewidth}p{0.07\linewidth}p{0.07\linewidth}p{0.09\linewidth}p{0.06\linewidth}p{0.21\linewidth}}
\toprule
ID & 领域 & 风格 & 参数完整性 & 期望初始 & 实际初始 & 期望最终 & 实际最终 & 缺参(初->终) & 结果 & 失败原因 \\
\midrule
MTM01 & medical & fuzzy & missing & 失败 & 失败 & 通过 & 通过 & 12->0 & 通过 & - \\
MTM02 & aerospace & precise & missing & 失败 & 失败 & 通过 & 通过 & 2->0 & 通过 & - \\
MTM03 & industrial\_ndt & precise & missing & 失败 & 失败 & 通过 & 通过 & 12->0 & 通过 & - \\
MTM04 & physics & precise & missing & 失败 & 失败 & 通过 & 通过 & 8->0 & 通过 & - \\
MTM05 & nuclear\_engineering & fuzzy & missing & 失败 & 失败 & 通过 & 通过 & 9->0 & 通过 & - \\
MTM06 & security\_screening & fuzzy & missing & 失败 & 失败 & 通过 & 通过 & 11->0 & 通过 & - \\
MTM07 & semiconductor & precise & missing & 失败 & 失败 & 通过 & 通过 & 3->0 & 通过 & - \\
MTM08 & education & fuzzy & missing & 失败 & 失败 & 通过 & 通过 & 11->0 & 通过 & - \\
MTM09 & environmental\_radiation & fuzzy & missing & 失败 & 失败 & 通过 & 通过 & 3->0 & 通过 & - \\
MTM10 & high\_energy\_physics & precise & missing & 失败 & 失败 & 通过 & 通过 & 10->0 & 通过 & - \\
MTM11 & materials\_update & precise & full & 通过 & 通过 & 通过 & 通过 & 0->0 & 通过 & - \\
MTM12 & overwrite\_reject & precise & full & 通过 & 通过 & 通过 & 通过 & 0->0 & 通过 & - \\
MTM13 & medical\_detector & fuzzy & missing & 失败 & 失败 & 通过 & 通过 & 13->0 & 通过 & - \\
MTM14 & industrial\_panel & precise & missing & 失败 & 失败 & 通过 & 通过 & 15->0 & 通过 & - \\
MTM15 & shielding\_stack & precise & missing & 失败 & 失败 & 通过 & 通过 & 14->0 & 通过 & - \\
MTM16 & nested\_target & precise & missing & 失败 & 失败 & 通过 & 通过 & 9->0 & 通过 & - \\
MTM17 & coaxial\_shield & fuzzy & missing & 失败 & 失败 & 通过 & 通过 & 14->0 & 通过 & - \\
MTM18 & boolean\_cavity & precise & missing & 失败 & 失败 & 通过 & 通过 & 14->0 & 通过 & - \\
MTM19 & ring\_complete & precise & full & 通过 & 通过 & 通过 & 通过 & 0->0 & 通过 & - \\
MTM20 & grid\_complete & precise & full & 通过 & 通过 & 通过 & 通过 & 0->0 & 通过 & - \\
MTM21 & layered\_detector & fuzzy & missing & 失败 & 失败 & 通过 & 通过 & 14->0 & 通过 & - \\
MTM22 & embedded\_sensor & precise & full & 通过 & 通过 & 通过 & 通过 & 0->0 & 通过 & - \\
MTM23 & layered\_shield\_complete & precise & full & 通过 & 通过 & 通过 & 通过 & 0->0 & 通过 & - \\
MTM24 & boolean\_complete & precise & full & 通过 & 通过 & 通过 & 通过 & 0->0 & 通过 & - \\
\bottomrule
\end{longtable}
\section{基准对比}
\begin{itemize}
\item pass\_rate=1.000
\item initial\_alignment\_rate=1.000
\item final\_alignment\_rate=1.000
\item avg\_turns\_per\_case=2.458
\item missing\_reduction\_rate=1.000
\item confirm\_apply\_success\_rate=1.000
\item rollback\_after\_confirm\_rate=0.000
\item internal\_leak\_turn\_rate=0.017
\item repeated\_assistant\_case\_rate=0.000
\item stale\_confirm\_case\_rate=0.000
\item delta\_pass\_rate=+0.050
\item regressions=无
\item improvements=无
\end{itemize}
\appendix
\section{附录：完整对话与内部处理链（自然语言）}
\subsection{案例 MTM01（medical）}
\textbf{对话：}
\begin{itemize}
\item 用户：我先要一个用于医学成像的初步辐照场景，细节还没定，你先告诉我还缺什么。
\item 系统：好的，我们来看看目前的进度。

**已完成的部分：**
*   我们已经确定了使用“盒体”作为几何结构类型。

**接下来还需要您提供的信息：**
1.  **几何**：盒体在x、y、z三个方向的具体尺寸。
2.  **材料**：这个盒体体积需要绑定哪种材料。
3.  **源**：您希望使用什么类型的放射源（例如伽马源、电子源等）。

请告诉我这些细节，我们就可以继续推进配置了。
\item 用户：采用 300 mm x 300 mm x 300 mm 的水箱，材料 G4\_WATER；源用束流 gamma，能量 6 MeV，位置 (0,0,-400) mm，方向 +z；物理列表 QBBC；输出 root。
\item 系统：我注意到这将会覆盖一些已经确认的配置。具体来说：
*   几何结构类型：从“未设定”改为“单一盒体”
*   盒体 x 方向尺寸：从“未设定”改为“300.0”

如果您确认应用这些修改，请直接回复“确认”；如果您希望保留原来的值，请回复“保持原值”。
\item 用户：确认
\item 系统：配置已完成。当前已确认的几何参数包括：几何结构类型、盒体在x、y、z三个方向的尺寸。
\end{itemize}
\textbf{内部链路（自然语言）：}
\begin{itemize}
\item 第 1 轮：动作=summarize\_progress，阶段=geometry，闭环状态=False。
\item 缺失字段：geometry.params.module\_x、geometry.params.module\_y、geometry.params.module\_z、materials.volume\_material\_map、source.type、source.particle、source.energy、source.position、source.direction、physics.physics\_list、output.format、output.path。
\item LLM 使用=True；回退=无；后端=llm\_semantic\_frame+runtime\_semantic。
\item 第 2 轮：动作=confirm\_overwrite，阶段=geometry，闭环状态=False。
\item 缺失字段：geometry.params.module\_x、geometry.params.module\_y、geometry.params.module\_z。
\item LLM 使用=True；回退=无；后端=llm\_slot\_frame+runtime\_semantic。
\item 第 3 轮：动作=finalize，阶段=materials，闭环状态=True。
\item 缺失字段：无。
\item LLM 使用=True；回退=无；后端=llm\_slot\_frame+runtime\_semantic+confirmed\_pending\_overwrite。
\end{itemize}
\subsection{案例 MTM02（aerospace）}
\textbf{对话：}
\begin{itemize}
\item 用户：I need a first-pass proton irradiation setup for a silicon device; keep it simple for now.
\item 系统：Geometry: updated geometry structure, box size along x, box size along y. Materials: updated selected materials, volume-to-material binding. Source: updated source type, particle type, source energy. Physics: updated physics list. Output: still needs output format, output path.
\item 用户：Use a silicon box 50 mm x 50 mm x 5 mm, beam source proton 100 MeV from (0,0,-100) mm along +z, physics FTFP\_BERT, output json.
\item 系统：You've requested a change that would overwrite a previously confirmed value.

Pending change:
- Box size along x: 1000.0 mm → 50.0 mm
- Box size along y: 1000.0 mm → 50.0 mm

Please reply **'confirm'** to apply this change, or **'keep existing'** to keep the current value.
\item 用户：confirm
\item 系统：Your configuration is complete. The following items are confirmed: box size along x, box size along y, box size along z, and the output path.
\end{itemize}
\textbf{内部链路（自然语言）：}
\begin{itemize}
\item 第 1 轮：动作=summarize\_progress，阶段=materials，闭环状态=False。
\item 缺失字段：output.format、output.path。
\item LLM 使用=True；回退=无；后端=llm\_slot\_frame+runtime\_semantic。
\item 第 2 轮：动作=confirm\_overwrite，阶段=source，闭环状态=False。
\item 缺失字段：无。
\item LLM 使用=True；回退=无；后端=llm\_slot\_frame+runtime\_semantic。
\item 第 3 轮：动作=finalize，阶段=physics，闭环状态=True。
\item 缺失字段：无。
\item LLM 使用=True；回退=无；后端=llm\_semantic\_frame+runtime\_semantic+confirmed\_pending\_overwrite。
\end{itemize}
\subsection{案例 MTM03（industrial\_ndt）}
\textbf{对话：}
\begin{itemize}
\item 用户：我要做一个 gamma 穿透铁靶的小测试，但我还没把完整参数说完。
\item 系统：好的，我们已确认了材料和粒子类型（gamma 射线和铁）。接下来还需要你补充几何结构的具体类型和尺寸，以及体积与材料的绑定关系。
\item 用户：请设为半径 40 mm、半长 80 mm 的铁圆柱，材料 G4\_Fe；点源 gamma，能量 2 MeV，位置 (0,0,-200) mm，方向 +z；物理列表 Shielding；输出 csv。
\item 系统：配置已完成。当前已确认的参数包括：体积与材料绑定、输出路径、输出格式和物理列表。
\end{itemize}
\textbf{内部链路（自然语言）：}
\begin{itemize}
\item 第 1 轮：动作=summarize\_progress，阶段=geometry，闭环状态=False。
\item 缺失字段：geometry.structure、geometry.params.module\_x、geometry.params.module\_y、geometry.params.module\_z、materials.volume\_material\_map、source.type、source.energy、source.position、source.direction、physics.physics\_list、output.format、output.path。
\item LLM 使用=True；回退=无；后端=llm\_slot\_frame+runtime\_semantic。
\item 第 2 轮：动作=finalize，阶段=materials，闭环状态=True。
\item 缺失字段：无。
\item LLM 使用=True；回退=无；后端=llm\_slot\_frame+runtime\_semantic。
\end{itemize}
\subsection{案例 MTM04（physics）}
\textbf{对话：}
\begin{itemize}
\item 用户：Gamma point source 1 MeV at (0,0,-100) toward +z; please build the rest.
\item 系统：Great. I've recorded your gamma source details: a 1 MeV point source at (0,0,-100) directed along +z.

**Progress Summary:**
*   **Source:** Complete (type, particle, energy, position, and direction are set).
*   **Geometry:** Needs the world volume shape (structure) and its X, Y, Z dimensions.
*   **Materials:** Needs to assign a material to that world volume.
*   **Physics \& Output:** Still require a physics list and output preferences.

To proceed, please specify the world volume's geometry and size. For example, you could say: "Use a box world volume with size 200 cm in each direction."
\item 用户：Set copper box 1 m x 1 m x 1 m, physics FTFP\_BERT, output hdf5.
\item 系统：Your configuration is complete. The following items are confirmed: the assignment of material to your geometry volume, the output file path and format, and the physics list.
\end{itemize}
\textbf{内部链路（自然语言）：}
\begin{itemize}
\item 第 1 轮：动作=summarize\_progress，阶段=geometry，闭环状态=False。
\item 缺失字段：geometry.structure、geometry.params.module\_x、geometry.params.module\_y、geometry.params.module\_z、materials.volume\_material\_map、physics.physics\_list、output.format、output.path。
\item LLM 使用=True；回退=无；后端=llm\_slot\_frame+runtime\_semantic。
\item 第 2 轮：动作=finalize，阶段=materials，闭环状态=True。
\item 缺失字段：无。
\item LLM 使用=True；回退=无；后端=llm\_slot\_frame+runtime\_semantic。
\end{itemize}
\subsection{案例 MTM05（nuclear\_engineering）}
\textbf{对话：}
\begin{itemize}
\item 用户：我想做一个中子屏蔽的概念验证，现在先给一个缺参开头。
\item 系统：目前我们已经确定了中子的粒子类型，并选择了盒体作为几何结构，材料也已绑定好，物理列表也同步完成。

接下来还需要你提供盒体在x、y、z三个方向的具体尺寸，以及源的详细类型和能量值。
\item 用户：使用 200 mm x 200 mm x 200 mm 的铅块，材料 G4\_Pb；点源 neutron，能量 14 MeV，位置 (0,0,-50) mm，方向 +z；物理列表 QGSP\_BERT\_HP；输出 xml。
\item 系统：我注意到这次设置会覆盖一些已经确认的参数。具体来说：

*   几何结构类型：从“未知”改为“单个盒体”
*   已选材料：从“G4\_Cu”改为“G4\_Pb”

如果你想应用这些修改，请直接回复“确认”；如果想保留原来的值，请直接说“保持原值”。
\item 用户：确认
\item 系统：配置已完成。当前已确认的参数包括：体积与材料绑定、盒体在x、y、z三个方向的尺寸。
\end{itemize}
\textbf{内部链路（自然语言）：}
\begin{itemize}
\item 第 1 轮：动作=summarize\_progress，阶段=geometry，闭环状态=False。
\item 缺失字段：geometry.params.module\_x、geometry.params.module\_y、geometry.params.module\_z、source.type、source.energy、source.position、source.direction、output.format、output.path。
\item LLM 使用=True；回退=无；后端=llm\_slot\_frame+runtime\_semantic。
\item 第 2 轮：动作=confirm\_overwrite，阶段=geometry，闭环状态=False。
\item 缺失字段：geometry.params.module\_x、geometry.params.module\_y、geometry.params.module\_z。
\item LLM 使用=True；回退=无；后端=llm\_slot\_frame+runtime\_semantic。
\item 第 3 轮：动作=finalize，阶段=materials，闭环状态=True。
\item 缺失字段：无。
\item LLM 使用=True；回退=无；后端=llm\_semantic\_frame+runtime\_semantic+confirmed\_pending\_overwrite。
\end{itemize}
\subsection{案例 MTM06（security\_screening）}
\textbf{对话：}
\begin{itemize}
\item 用户：Need a quick cargo-screening toy setup with forward photons, but I have not fixed the target shape, material, source energy, or output format yet.
\item 系统：Geometry: updated geometry structure; still needs box size along x, box size along y. Materials: still needs volume-to-material binding. Source: updated particle type; still needs source type, source energy. Physics: still needs physics list. Output: still needs output format, output path.
\item 用户：Use steel cylinder target radius 60 mm half-length 120 mm, beam gamma 3 MeV from (0,0,-500) mm to +z, physics Shielding, output root.
\item 系统：This would overwrite an already confirmed value. Pending change: geometry structure: unknown -> single\_tubs; box size along x: None -> 1000.0. Reply 'confirm' to apply it, or say 'keep existing' to keep the current value.
\item 用户：confirm
\item 系统：Configuration complete. Confirmed items: volume-to-material binding, half length, radius, and box size along z.
\end{itemize}
\textbf{内部链路（自然语言）：}
\begin{itemize}
\item 第 1 轮：动作=summarize\_progress，阶段=geometry，闭环状态=False。
\item 缺失字段：geometry.params.module\_x、geometry.params.module\_y、geometry.params.module\_z、materials.volume\_material\_map、source.type、source.energy、source.position、source.direction、physics.physics\_list、output.format、output.path。
\item LLM 使用=True；回退=无；后端=llm\_slot\_frame+runtime\_semantic。
\item 第 2 轮：动作=confirm\_overwrite，阶段=geometry，闭环状态=False。
\item 缺失字段：geometry.params.module\_x、geometry.params.module\_y、geometry.params.module\_z。
\item LLM 使用=True；回退=无；后端=llm\_slot\_frame+runtime\_semantic。
\item 第 3 轮：动作=finalize，阶段=materials，闭环状态=True。
\item 缺失字段：无。
\item LLM 使用=True；回退=无；后端=llm\_semantic\_frame+runtime\_semantic+confirmed\_pending\_overwrite。
\end{itemize}
\subsection{案例 MTM07（semiconductor）}
\textbf{对话：}
\begin{itemize}
\item 用户：我想搭一个硅薄片的中子照射测试，但先只给一个简述。
\item 系统：目前我们已经确定了几何结构（盒体尺寸）、材料（硅）和源设置（中子源参数）。接下来还需要你补充物理列表、输出格式和输出路径这三项信息。
\item 用户：使用 20 mm x 20 mm x 1 mm 的硅片，材料 G4\_Si；各向同性 neutron，能量 2 MeV，位置 (0,0,0) mm；物理列表 QGSP\_BERT\_HP；输出 json。
\item 系统：我注意到这些修改会覆盖已确认的数值，具体如下：

- 盒体 x 方向尺寸：从 1000.0 改为 20.0
- 盒体 y 方向尺寸：从 1000.0 改为 20.0

如果你确认应用这些修改，请直接回复“确认”；如果想保留原来的数值，请直接说“保持原值”。
\item 用户：确认
\item 系统：配置已完成。当前已确认的参数包括：源位置、源能量、源类型以及盒体 z 方向尺寸。
\end{itemize}
\textbf{内部链路（自然语言）：}
\begin{itemize}
\item 第 1 轮：动作=summarize\_progress，阶段=materials，闭环状态=False。
\item 缺失字段：physics.physics\_list、output.format、output.path。
\item LLM 使用=True；回退=无；后端=llm\_slot\_frame+runtime\_semantic。
\item 第 2 轮：动作=confirm\_overwrite，阶段=source，闭环状态=False。
\item 缺失字段：无。
\item LLM 使用=True；回退=无；后端=llm\_slot\_frame+runtime\_semantic。
\item 第 3 轮：动作=finalize，阶段=physics，闭环状态=True。
\item 缺失字段：无。
\item LLM 使用=True；回退=无；后端=llm\_semantic\_frame+runtime\_semantic+confirmed\_pending\_overwrite。
\end{itemize}
\subsection{案例 MTM08（education）}
\textbf{对话：}
\begin{itemize}
\item 用户：For classroom demo, I want gamma attenuation in copper; you can ask only essential missing parts.
\item 系统：Here's a summary of your setup so far and what's still needed:

**Completed:**
*   **Materials:** Copper is selected.
*   **Source:** The particle type is set to gamma.
*   **Physics:** A physics list has been selected.

**Still needed to configure:**
*   **Geometry:** The structure type and the box size along the x, y, and z axes.
*   **Materials:** You'll need to assign the copper material to a specific volume.
*   **Source:** The source type, energy, and position.
*   **Output:** The format and file path for your results.
\item 用户：Okay, copper box 100 mm x 100 mm x 100 mm, point gamma 1 MeV at (0,0,-200) mm, direction +z, physics FTFP\_BERT, output csv.
\item 系统：You're about to change a confirmed setting. The physics list would be updated from "Shielding" to "FTFP\_BERT". Please reply 'confirm' to apply this change, or 'keep existing' to keep the current value.
\item 用户：confirm
\item 系统：Your configuration is complete. The physics list, output path, output format, and source direction have been confirmed.
\end{itemize}
\textbf{内部链路（自然语言）：}
\begin{itemize}
\item 第 1 轮：动作=summarize\_progress，阶段=geometry，闭环状态=False。
\item 缺失字段：geometry.structure、geometry.params.module\_x、geometry.params.module\_y、geometry.params.module\_z、materials.volume\_material\_map、source.type、source.energy、source.position、source.direction、output.format、output.path。
\item LLM 使用=True；回退=无；后端=llm\_slot\_frame+runtime\_semantic。
\item 第 2 轮：动作=confirm\_overwrite，阶段=materials，闭环状态=False。
\item 缺失字段：无。
\item LLM 使用=True；回退=无；后端=llm\_slot\_frame+runtime\_semantic。
\item 第 3 轮：动作=finalize，阶段=source，闭环状态=True。
\item 缺失字段：无。
\item LLM 使用=True；回退=无；后端=llm\_semantic\_frame+runtime\_semantic+confirmed\_pending\_overwrite。
\end{itemize}
\subsection{案例 MTM09（environmental\_radiation）}
\textbf{对话：}
\begin{itemize}
\item 用户：我想先搭一个空气体积中的低能 gamma 传播场景，细节稍后补。
\item 系统：目前我们已经完成了几何、材料和源的设置，接下来还需要你补充物理列表、输出格式和输出路径这三项信息。
\item 用户：设为 500 mm x 500 mm x 500 mm 的空气盒，材料 G4\_AIR；点源 gamma，能量 0.5 MeV，位置 (0,0,-300) mm，方向 +z；物理列表 QBBC；输出 root。
\item 系统：我注意到这会覆盖已确认的设置。具体修改如下：

- 盒体 x 方向尺寸：1000.0 → 500.0
- 盒体 y 方向尺寸：1000.0 → 500.0

请确认是否应用这些修改：
• 如果确认修改，请直接回复“确认”。
• 如果想保留原值，请直接回复“保持原值”。
\item 用户：确认
\item 系统：好的，配置已完成！当前已确认的设置包括：源位置、源能量、以及盒体的 z、y、x 方向尺寸。
\end{itemize}
\textbf{内部链路（自然语言）：}
\begin{itemize}
\item 第 1 轮：动作=summarize\_progress，阶段=materials，闭环状态=False。
\item 缺失字段：physics.physics\_list、output.format、output.path。
\item LLM 使用=True；回退=无；后端=llm\_slot\_frame+runtime\_semantic。
\item 第 2 轮：动作=confirm\_overwrite，阶段=source，闭环状态=False。
\item 缺失字段：无。
\item LLM 使用=True；回退=无；后端=llm\_slot\_frame+runtime\_semantic。
\item 第 3 轮：动作=finalize，阶段=physics，闭环状态=True。
\item 缺失字段：无。
\item LLM 使用=True；回退=无；后端=llm\_semantic\_frame+runtime\_semantic+confirmed\_pending\_overwrite。
\end{itemize}
\subsection{案例 MTM10（high\_energy\_physics）}
\textbf{对话：}
\begin{itemize}
\item 用户：Set geometry as a sphere target; I have not decided radius or source details yet.
\item 系统：Here's where we stand:

**Completed this step:**
*   **Geometry:** You've set the target to a sphere.

**Still needed to proceed:**
*   **Geometry:** The sphere's radius.
*   **Materials:** Which material to assign to the sphere.
*   **Source:** The source type, particle type, and its energy.
*   **Physics:** A physics list.
*   **Output:** The output format and file path.
\item 用户：Sphere target radius 80 mm, lead material, electron point source 5 MeV at (0,0,-200) mm toward +z, physics FTFP\_BERT, output hdf5.
\item 系统：Your configuration is complete. The following items have been confirmed: the material binding for your volumes, the output file path and format, and the physics list.
\end{itemize}
\textbf{内部链路（自然语言）：}
\begin{itemize}
\item 第 1 轮：动作=summarize\_progress，阶段=geometry，闭环状态=False。
\item 缺失字段：geometry.params.child\_rmax、materials.volume\_material\_map、source.type、source.particle、source.energy、source.position、source.direction、physics.physics\_list、output.format、output.path。
\item LLM 使用=True；回退=无；后端=llm\_slot\_frame+runtime\_semantic。
\item 第 2 轮：动作=finalize，阶段=materials，闭环状态=True。
\item 缺失字段：无。
\item LLM 使用=True；回退=无；后端=llm\_slot\_frame+runtime\_semantic。
\end{itemize}
\subsection{案例 MTM11（materials\_update）}
\textbf{对话：}
\begin{itemize}
\item 用户：Create a complete setup: 100 mm x 100 mm x 100 mm copper box, point gamma 1 MeV at (0,0,-200) mm toward +z, physics QBBC, output json.
\item 系统：Current committed config is: physics list FTFP\_BERT, output format json, source type point, particle gamma. If you want changes, please specify which field to update.
\item 用户：Change the target material to G4\_Al and change the output format to root.
\item 系统：You're about to overwrite two confirmed settings:

*   **Selected materials:** from copper (G4\_Cu) to aluminum (G4\_Al)
*   **Output format:** from JSON to ROOT

Please reply **"confirm"** to apply these changes, or **"keep existing"** to keep the current values.
\item 用户：confirm
\item 系统：Your configuration is complete. The following items are confirmed: the volume-to-material binding, output path, output format, and selected materials.
\end{itemize}
\textbf{内部链路（自然语言）：}
\begin{itemize}
\item 第 1 轮：动作=finalize，阶段=materials，闭环状态=True。
\item 缺失字段：无。
\item LLM 使用=True；回退=无；后端=llm\_slot\_frame+runtime\_semantic。
\item 第 2 轮：动作=confirm\_overwrite，阶段=source，闭环状态=False。
\item 缺失字段：无。
\item LLM 使用=True；回退=无；后端=llm\_slot\_frame+runtime\_semantic。
\item 第 3 轮：动作=finalize，阶段=physics，闭环状态=True。
\item 缺失字段：无。
\item LLM 使用=True；回退=无；后端=llm\_semantic\_frame+runtime\_semantic+confirmed\_pending\_overwrite。
\end{itemize}
\subsection{案例 MTM12（overwrite\_reject）}
\textbf{对话：}
\begin{itemize}
\item 用户：请创建一个完整配置：100 mm x 100 mm x 100 mm 的铜盒，材料 G4\_Cu；点源 gamma，能量 1 MeV，位置 (0,0,-200) mm，方向 +z；物理列表 QBBC；输出 json。
\item 系统：配置已完成。当前已确认的设置包括：体积与材料绑定、输出路径和输出格式。
\item 用户：把材料改成 G4\_Al，并把输出格式改成 root。
\item 系统：我注意到这次修改会覆盖已经确认的设置。具体变动如下：

- 已选材料：从 `G4\_Cu` 改为 `G4\_Al`
- 输出格式：从 `json` 改为 `root`

请确认您的选择：
- 如果同意应用以上修改，请回复“确认”。
- 如果希望保留原来的值，请回复“保持原值”。
\item 用户：保持原值
\item 系统：好的，已保留“已选材料”和“输出格式”的原有值，本次修改未被应用。
\end{itemize}
\textbf{内部链路（自然语言）：}
\begin{itemize}
\item 第 1 轮：动作=finalize，阶段=materials，闭环状态=True。
\item 缺失字段：无。
\item LLM 使用=True；回退=无；后端=llm\_slot\_frame+runtime\_semantic。
\item 第 2 轮：动作=confirm\_overwrite，阶段=source，闭环状态=False。
\item 缺失字段：geometry.ask.boolean.solid\_a\_size、geometry.ask.boolean.solid\_b\_size。
\item LLM 使用=True；回退=无；后端=llm\_slot\_frame+runtime\_semantic。
\item 第 3 轮：动作=reject\_overwrite，阶段=physics，闭环状态=True。
\item 缺失字段：无。
\item LLM 使用=True；回退=无；后端=llm\_semantic\_frame+runtime\_semantic+rejected\_pending\_overwrite。
\end{itemize}
\subsection{案例 MTM13（medical\_detector）}
\textbf{对话：}
\begin{itemize}
\item 用户：I want a ring of detector modules around a small source region, but I have not fixed the ring radius or module size yet.
\item 系统：I've updated your geometry to a ring structure. To continue, I need:

*   **Geometry:** The size of each detector module and how many will be in the ring.
*   **Materials:** To assign a material to your geometry volume.
*   **Source:** The type of source, the particle to simulate, and its energy.
*   **Physics \& Output:** A physics list, plus an output format and file path.
\item 用户：Use a ring of 12 modules, each 8 mm x 10 mm x 2 mm, radius 40 mm, clearance 1 mm, material G4\_WATER; point gamma source 0.662 MeV at (0,0,0) mm; physics QBBC; output json.
\item 系统：You've requested a change that would overwrite a previously confirmed geometry.

**Current geometry:**
*   Module: Box (7.83 x 4.72 x 6.86 mm)
*   Ring: 4 modules, radius 67.49 mm, clearance 0.19 mm

**New geometry:**
*   Module: Box (8 x 10 x 2 mm)
*   Ring: 12 modules, radius 40 mm, clearance 1 mm

Please reply **'confirm'** to apply this change, or **'keep existing'** to keep the current values.
\item 用户：confirm
\item 系统：Your configuration is complete. The following items are confirmed: ring clearance, ring radius, ring module count, and ring module size.
\end{itemize}
\textbf{内部链路（自然语言）：}
\begin{itemize}
\item 第 1 轮：动作=summarize\_progress，阶段=materials，闭环状态=False。
\item 缺失字段：materials.volume\_material\_map、source.type、source.particle、source.energy、source.position、source.direction、physics.physics\_list、output.format、output.path、geometry.ask.ring.module\_size、geometry.ask.ring.count、geometry.ask.ring.radius、geometry.ask.ring.clearance。
\item LLM 使用=True；回退=无；后端=llm\_slot\_frame+runtime\_semantic。
\item 第 2 轮：动作=confirm\_overwrite，阶段=source，闭环状态=False。
\item 缺失字段：无。
\item LLM 使用=True；回退=无；后端=llm\_slot\_frame+runtime\_semantic。
\item 第 3 轮：动作=finalize，阶段=physics，闭环状态=True。
\item 缺失字段：无。
\item LLM 使用=True；回退=无；后端=llm\_semantic\_frame+runtime\_semantic+confirmed\_pending\_overwrite。
\end{itemize}
\subsection{案例 MTM14（industrial\_panel）}
\textbf{对话：}
\begin{itemize}
\item 用户：我想做一个二维阵列探测板，但 pitch 和材料还没定。
\item 系统：好的，目前进度如下：

**已确认的部分**：
*   **几何**：结构类型已确定为二维阵列探测板。盒体在x、y、z三个方向上的尺寸也已确认。

**接下来需要您补充的信息**：
*   **几何**：阵列模块的具体尺寸、x和y方向的阵列数量。
*   **材料**：体积与材料的绑定关系。
*   **源**：源的类型、粒子种类和能量。
*   **物理**：物理过程列表。
*   **输出**：输出文件的格式和保存路径。

请随时告诉我这些信息，我们继续配置。
\item 用户：设成 4 x 3 的盒体阵列，每个模块 10 mm x 10 mm x 2 mm，pitch\_x 12 mm，pitch\_y 14 mm，间隙 1 mm，材料 G4\_Si；点源 gamma 0.511 MeV，位置 (0,0,-80) mm，方向 +z；物理列表 QBBC；输出 root。
\item 系统：我注意到这次修改会覆盖之前已经确认的几何参数。具体来说，阵列模块的尺寸、阵列数量、pitch 和间隙都会被更新。

如果您确认要应用这些新值，请直接回复“确认”；如果您希望保留原来的设置，请回复“保持原值”。
\item 用户：确认
\item 系统：好的，配置已完成！当前已确认的参数包括：阵列间隙、y方向pitch、x方向pitch、y方向阵列数量和x方向阵列数量。
\end{itemize}
\textbf{内部链路（自然语言）：}
\begin{itemize}
\item 第 1 轮：动作=summarize\_progress，阶段=materials，闭环状态=False。
\item 缺失字段：materials.volume\_material\_map、source.type、source.particle、source.energy、source.position、source.direction、physics.physics\_list、output.format、output.path、geometry.ask.grid.module\_size、geometry.ask.grid.count\_x、geometry.ask.grid.count\_y、geometry.ask.grid.pitch\_x、geometry.ask.grid.pitch\_y、geometry.ask.grid.clearance。
\item LLM 使用=True；回退=无；后端=llm\_semantic\_frame+runtime\_semantic。
\item 第 2 轮：动作=confirm\_overwrite，阶段=source，闭环状态=False。
\item 缺失字段：无。
\item LLM 使用=True；回退=无；后端=llm\_slot\_frame+runtime\_semantic。
\item 第 3 轮：动作=finalize，阶段=physics，闭环状态=True。
\item 缺失字段：无。
\item LLM 使用=True；回退=无；后端=llm\_semantic\_frame+runtime\_semantic+confirmed\_pending\_overwrite。
\end{itemize}
\subsection{案例 MTM15（shielding\_stack）}
\textbf{对话：}
\begin{itemize}
\item 用户：Build a stacked absorber sandwich along z; I will provide the layer sizes next.
\item 系统：You've set up a stacked absorber sandwich geometry along the z-axis. Here's what's complete and what's still needed:

**Completed:**
*   **Geometry:** The overall structure is set. You've also confirmed the outer box dimensions.

**Still Needed:**
*   **Geometry:** The footprint and thicknesses for each stack layer.
*   **Materials:** Which material to assign to each volume.
*   **Source:** The source type, particle type, and its energy.
*   **Physics:** The physics list to use.
*   **Output:** The format and file path for your results.
\item 用户：Use stacked layers with footprint 20 mm x 20 mm, thicknesses 2 mm, 4 mm, 6 mm, layer clearance 0.5 mm, outer box 30 mm x 30 mm x 20 mm, container clearance 1 mm, material G4\_Al; beam gamma 1.25 MeV from (0,0,-120) mm along +z; physics Shielding; output csv.
\item 系统：This would overwrite an already confirmed geometry configuration.

**Current values:**
*   Stack footprint: \textasciitilde{}16.33 mm x \textasciitilde{}10.28 mm
*   Layer thicknesses: \textasciitilde{}6.68 mm, \textasciitilde{}1.19 mm, \textasciitilde{}5.59 mm
*   Layer clearance: \textasciitilde{}0.73 mm
*   Outer box: \textasciitilde{}18.07 mm x \textasciitilde{}45.52 mm x \textasciitilde{}16.66 mm
*   Container clearance: \textasciitilde{}0.87 mm

**New values:**
*   Stack footprint: 20.0 mm x 20.0 mm
*   Layer thicknesses: 2.0 mm, 4.0 mm, 6.0 mm
*   Layer clearance: 0.5 mm
*   Outer box: 30.0 mm x 30.0 mm x 20.0 mm
*   Container clearance: 1.0 mm

Reply **'confirm'** to apply these new values, or **'keep existing'** to keep the current geometry.
\item 用户：confirm
\item 系统：Your configuration is complete. The following geometry items have been confirmed: stack container clearance, stack container size, stack layer clearance, and stack layer thicknesses.
\end{itemize}
\textbf{内部链路（自然语言）：}
\begin{itemize}
\item 第 1 轮：动作=summarize\_progress，阶段=materials，闭环状态=False。
\item 缺失字段：materials.volume\_material\_map、source.type、source.particle、source.energy、source.position、source.direction、physics.physics\_list、output.format、output.path、geometry.ask.stack.footprint、geometry.ask.stack.thicknesses、geometry.ask.stack.layer\_clearance、geometry.ask.stack.parent\_size、geometry.ask.stack.nest\_clearance。
\item LLM 使用=True；回退=无；后端=llm\_slot\_frame+runtime\_semantic。
\item 第 2 轮：动作=confirm\_overwrite，阶段=source，闭环状态=False。
\item 缺失字段：无。
\item LLM 使用=True；回退=无；后端=llm\_slot\_frame+runtime\_semantic。
\item 第 3 轮：动作=finalize，阶段=physics，闭环状态=True。
\item 缺失字段：无。
\item LLM 使用=True；回退=无；后端=llm\_semantic\_frame+runtime\_semantic+confirmed\_pending\_overwrite。
\end{itemize}
\subsection{案例 MTM16（nested\_target）}
\textbf{对话：}
\begin{itemize}
\item 用户：我想在一个盒子里嵌一个圆柱靶，先别急着补全。
\item 系统：目前已完成几何结构类型、源能量、位置和方向的设置。接下来还需要补充体积与材料绑定、粒子类型、物理列表、输出格式和输出路径这些信息。
\item 用户：外盒 120 mm x 120 mm x 120 mm，内嵌铜圆柱半径 20 mm、半长 40 mm，嵌套间隙 2 mm；材料 G4\_Cu；束流 gamma 1 MeV，位置 (0,0,-150) mm，方向 +z；物理列表 QBBC；输出 json。
\item 系统：我注意到这可能会覆盖已经确认的设置。具体修改如下：

- 源类型：从“点源”改为“束流”
- 源位置：从 (0,0,-100) mm 改为 (0,0,-150) mm

如果你确认应用这些修改，请直接回复“确认”；如果想保留原来的值，请回复“保持原值”。
\item 用户：确认
\item 系统：配置已完成。当前已确认的设置包括：源位置、源类型、嵌套间隙和内嵌圆柱半长。
\end{itemize}
\textbf{内部链路（自然语言）：}
\begin{itemize}
\item 第 1 轮：动作=summarize\_progress，阶段=materials，闭环状态=False。
\item 缺失字段：materials.volume\_material\_map、source.particle、physics.physics\_list、output.format、output.path、geometry.ask.nest.parent\_size、geometry.ask.nest.child\_radius、geometry.ask.nest.child\_half\_length、geometry.ask.nest.clearance。
\item LLM 使用=True；回退=无；后端=llm\_slot\_frame+runtime\_semantic。
\item 第 2 轮：动作=confirm\_overwrite，阶段=source，闭环状态=False。
\item 缺失字段：无。
\item LLM 使用=True；回退=无；后端=llm\_slot\_frame+runtime\_semantic。
\item 第 3 轮：动作=finalize，阶段=physics，闭环状态=True。
\item 缺失字段：无。
\item LLM 使用=True；回退=无；后端=llm\_semantic\_frame+runtime\_semantic+confirmed\_pending\_overwrite。
\end{itemize}
\subsection{案例 MTM17（coaxial\_shield）}
\textbf{对话：}
\begin{itemize}
\item 用户：I need a shell-style shielding toy model, but the shell thicknesses are still open.
\item 系统：Here's what we have so far:

**Geometry:** We've set up a shell structure. Next, we need the shell's inner radius and its thicknesses.

**Materials:** We still need to assign materials to each volume.

**Source:** We still need the source type, particle type, and its energy and position.

**Physics:** The physics list is now configured.

**Output:** We still need to choose an output format and a file path.
\item 用户：Use a shell with inner radius 15 mm, thicknesses 5 mm and 8 mm, half length 40 mm, nested core radius 10 mm, nested core half length 30 mm, clearance 1 mm, material G4\_Pb; isotropic neutron source 2 MeV at (0,0,0) mm; physics Shielding; output hdf5.
\item 系统：You're about to overwrite a geometry that was already set.

**Current geometry:**
*   **Shell:** Inner radius 7.15 mm, thicknesses [1.18, 3.43, 0.86] mm, half-length 22.36 mm.
*   **Nested Core:** Radius 11.60 mm, half-length 2.68 mm.
*   **Clearance:** 1.01 mm.

**New geometry:**
*   **Shell:** Inner radius 15.0 mm, thicknesses [5.0, 8.0, 0.83] mm, half-length 40.0 mm.
*   **Nested Core:** Radius 10.0 mm, half-length 30.0 mm.
*   **Clearance:** 1.0 mm.

Please reply **'confirm'** to apply this change, or **'keep existing'** to keep the current geometry.
\item 用户：confirm
\item 系统：Configuration complete. The following items are confirmed: shell clearance, nested core half length, nested core radius, and shell half length.
\end{itemize}
\textbf{内部链路（自然语言）：}
\begin{itemize}
\item 第 1 轮：动作=summarize\_progress，阶段=materials，闭环状态=False。
\item 缺失字段：materials.volume\_material\_map、source.type、source.particle、source.energy、source.position、source.direction、output.format、output.path、geometry.ask.shell.inner\_radius、geometry.ask.shell.thicknesses、geometry.ask.shell.half\_length、geometry.ask.shell.child\_radius、geometry.ask.shell.child\_half\_length、geometry.ask.shell.clearance。
\item LLM 使用=True；回退=无；后端=llm\_slot\_frame+runtime\_semantic。
\item 第 2 轮：动作=confirm\_overwrite，阶段=source，闭环状态=False。
\item 缺失字段：无。
\item LLM 使用=True；回退=无；后端=llm\_slot\_frame+runtime\_semantic。
\item 第 3 轮：动作=finalize，阶段=physics，闭环状态=True。
\item 缺失字段：无。
\item LLM 使用=True；回退=无；后端=llm\_semantic\_frame+runtime\_semantic+confirmed\_pending\_overwrite。
\end{itemize}
\subsection{案例 MTM18（boolean\_cavity）}
\textbf{对话：}
\begin{itemize}
\item 用户：我想做一个挖空的盒体几何，但先只给你意图。
\item 系统：好的，我们已同步了您的几何结构类型、源位置、源方向和源类型。

接下来，为了完成挖空盒体的几何定义，还需要您补充 boolean solid A 的 x 尺寸和 y 尺寸。
\item 用户：做减法布尔体：A 为 40 mm x 40 mm x 40 mm 盒体，B 为 20 mm x 20 mm x 20 mm 盒体，从 A 中减去 B；材料 G4\_Al；点源 gamma 0.3 MeV，位置 (0,0,-60) mm，方向 +z；物理列表 QBBC；输出 xml。
\item 系统：我注意到这可能会覆盖已经确认的设置。具体来说：

1. **源位置**：将从 (0, 0, -100) mm 改为 (0, 0, -60) mm
2. **几何结构**：两个盒体的尺寸将更新为您指定的新数值

如果您确认应用这些修改，请直接回复“确认”；如果希望保留原来的值，请回复“保持原值”。
\item 用户：确认
\item 系统：配置已完成！当前已确认的设置包括：源位置、布尔 solid B 尺寸、布尔 solid A 尺寸以及输出路径。
\end{itemize}
\textbf{内部链路（自然语言）：}
\begin{itemize}
\item 第 1 轮：动作=summarize\_progress，阶段=geometry，闭环状态=False。
\item 缺失字段：geometry.params.bool\_a\_x、geometry.params.bool\_a\_y、geometry.params.bool\_a\_z、geometry.params.bool\_b\_x、geometry.params.bool\_b\_y、geometry.params.bool\_b\_z、materials.volume\_material\_map、source.particle、source.energy、physics.physics\_list、output.format、output.path、geometry.ask.boolean.solid\_a\_size、geometry.ask.boolean.solid\_b\_size。
\item LLM 使用=True；回退=无；后端=llm\_slot\_frame+runtime\_semantic。
\item 第 2 轮：动作=confirm\_overwrite，阶段=materials，闭环状态=False。
\item 缺失字段：无。
\item LLM 使用=True；回退=无；后端=llm\_slot\_frame+runtime\_semantic。
\item 第 3 轮：动作=finalize，阶段=source，闭环状态=True。
\item 缺失字段：无。
\item LLM 使用=True；回退=无；后端=llm\_semantic\_frame+runtime\_semantic+confirmed\_pending\_overwrite。
\end{itemize}
\subsection{案例 MTM19（ring\_complete）}
\textbf{对话：}
\begin{itemize}
\item 用户：Create a complete ring setup: 16 modules, each 6 mm x 6 mm x 2 mm, radius 32 mm, clearance 0.5 mm, material G4\_Si; point gamma source 0.511 MeV at (0,0,0) mm; physics QBBC; output json.
\item 系统：Your configuration is complete. The following items have been confirmed: volume-to-material binding, output path, clearance, and radius.
\end{itemize}
\textbf{内部链路（自然语言）：}
\begin{itemize}
\item 第 1 轮：动作=finalize，阶段=materials，闭环状态=True。
\item 缺失字段：无。
\item LLM 使用=True；回退=无；后端=llm\_slot\_frame+runtime\_semantic。
\end{itemize}
\subsection{案例 MTM20（grid\_complete）}
\textbf{对话：}
\begin{itemize}
\item 用户：请创建完整的阵列配置：3 x 3 盒体阵列，每个模块 12 mm x 12 mm x 3 mm，pitch\_x 15 mm，pitch\_y 15 mm，间隙 1 mm，材料 G4\_CsI；点源 gamma 0.662 MeV，位置 (0,0,-90) mm，方向 +z；物理列表 QBBC；输出 root。
\item 系统：配置已完成。当前已确认：体积与材料绑定、输出路径、clearance、y 方向 pitch。
\end{itemize}
\textbf{内部链路（自然语言）：}
\begin{itemize}
\item 第 1 轮：动作=finalize，阶段=materials，闭环状态=True。
\item 缺失字段：无。
\item LLM 使用=True；回退=无；后端=llm\_slot\_frame+runtime\_semantic。
\end{itemize}
\subsection{案例 MTM21（layered\_detector）}
\textbf{对话：}
\begin{itemize}
\item 用户：I want a stacked detector slab, but the individual layer thicknesses are not fixed yet.
\item 系统：You've set up a stacked detector slab geometry. Next, we need to define the stack's footprint and layer thicknesses, assign materials to the volumes, and configure the particle source, physics list, and output settings.
\item 用户：Use stack footprint 25 mm x 25 mm, thicknesses 1 mm, 2 mm, 1 mm, layer clearance 0.2 mm, outer box 30 mm x 30 mm x 12 mm, container clearance 0.5 mm, material G4\_Si; beam electron 5 MeV from (0,0,-50) mm along +z; physics FTFP\_BERT; output json.
\item 系统：You're about to overwrite a geometry you already confirmed.

The pending change would update your stack and container dimensions:
*   Stack footprint: 16.33 mm x 10.28 mm → 25.0 mm x 25.0 mm
*   Stack layer thicknesses: [6.68, 1.19, 5.59] mm → [1.0, 2.0, 1.0] mm
*   Stack layer clearance: 0.73 mm → 0.2 mm
*   Outer container size: 18.07 mm x 45.52 mm x 16.66 mm → 30.0 mm x 30.0 mm x 12.0 mm
*   Container clearance: 0.87 mm → 0.5 mm

Please reply **'confirm'** to apply this change, or **'keep existing'** to keep the current geometry.
\item 用户：confirm
\item 系统：Your configuration is complete. The following geometry settings are confirmed: stack container clearance, stack container size, stack layer clearance, and stack layer thicknesses.
\end{itemize}
\textbf{内部链路（自然语言）：}
\begin{itemize}
\item 第 1 轮：动作=summarize\_progress，阶段=materials，闭环状态=False。
\item 缺失字段：materials.volume\_material\_map、source.type、source.particle、source.energy、source.position、source.direction、physics.physics\_list、output.format、output.path、geometry.ask.stack.footprint、geometry.ask.stack.thicknesses、geometry.ask.stack.layer\_clearance、geometry.ask.stack.parent\_size、geometry.ask.stack.nest\_clearance。
\item LLM 使用=True；回退=无；后端=llm\_semantic\_frame+runtime\_semantic。
\item 第 2 轮：动作=confirm\_overwrite，阶段=source，闭环状态=False。
\item 缺失字段：无。
\item LLM 使用=True；回退=无；后端=llm\_slot\_frame+runtime\_semantic。
\item 第 3 轮：动作=finalize，阶段=physics，闭环状态=True。
\item 缺失字段：无。
\item LLM 使用=True；回退=无；后端=llm\_semantic\_frame+runtime\_semantic+confirmed\_pending\_overwrite。
\end{itemize}
\subsection{案例 MTM22（embedded\_sensor）}
\textbf{对话：}
\begin{itemize}
\item 用户：请创建完整的嵌套配置：外盒 80 mm x 80 mm x 80 mm，内嵌铅圆柱半径 15 mm、半长 25 mm，嵌套间隙 1 mm；材料 G4\_Pb；点源 gamma 1 MeV，位置 (0,0,-100) mm，方向 +z；物理列表 Shielding；输出 csv。
\item 系统：配置已完成。当前已确认的设置包括：体积与材料绑定、输出路径、嵌套间隙（clearance）以及圆柱半径。
\end{itemize}
\textbf{内部链路（自然语言）：}
\begin{itemize}
\item 第 1 轮：动作=finalize，阶段=materials，闭环状态=True。
\item 缺失字段：无。
\item LLM 使用=True；回退=无；后端=llm\_slot\_frame+runtime\_semantic。
\end{itemize}
\subsection{案例 MTM23（layered\_shield\_complete）}
\textbf{对话：}
\begin{itemize}
\item 用户：Create a complete shell configuration: inner radius 12 mm, thicknesses 4 mm and 6 mm, half length 35 mm, nested core radius 8 mm, nested core half length 25 mm, clearance 1 mm, material G4\_Fe; isotropic gamma source 0.8 MeV at (0,0,0) mm; physics QBBC; output hdf5.
\item 系统：Your configuration is complete. The following items are confirmed: box size along z, box size along y, box size along x, and volume-to-material binding.
\end{itemize}
\textbf{内部链路（自然语言）：}
\begin{itemize}
\item 第 1 轮：动作=finalize，阶段=materials，闭环状态=True。
\item 缺失字段：无。
\item LLM 使用=True；回退=无；后端=llm\_slot\_frame+runtime\_semantic。
\end{itemize}
\subsection{案例 MTM24（boolean\_complete）}
\textbf{对话：}
\begin{itemize}
\item 用户：请创建完整布尔配置：做并集布尔体，A 为 30 mm x 30 mm x 20 mm 盒体，B 为 10 mm x 10 mm x 40 mm 盒体，与 A 做 union；材料 G4\_Al；束流 gamma 0.5 MeV，位置 (0,0,-70) mm，方向 +z；物理列表 QBBC；输出 json。
\item 系统：配置已完成。当前已确认的参数包括：盒体在 x、y、z 方向的尺寸，以及体积与材料的绑定关系。
\end{itemize}
\textbf{内部链路（自然语言）：}
\begin{itemize}
\item 第 1 轮：动作=finalize，阶段=materials，闭环状态=True。
\item 缺失字段：无。
\item LLM 使用=True；回退=无；后端=llm\_slot\_frame+runtime\_semantic。
\end{itemize}
\end{document}